% paper1_methods.tex -- Materials and Methods of Paper 1, input by
% papers/paper1_indicators.tex.
\section{Materials and Methods}

\subsection{Trial and sampling}
The indicators were measured in a trial comparing three management systems at
each of the three sites of the trial network: a conventional system with
tillage and mineral fertiliser only (CONV), a conservation system with reduced
tillage and retained residues (CONS), and an integrated system combining
reduced tillage, residues, organic inputs and fertiliser (INTG). The systems
were arranged in randomised complete blocks with four blocks per site. Topsoil
was sampled to a depth of \SI{0.20}{\metre} from every plot. The data are held
in \texttt{data/trials/t1\_indicators.tsv}.

\subsection{Selection of indicators}
Candidate indicators were retained only if they met four criteria: they respond
to management within a few seasons, they relate to at least one soil function,
they can be measured at moderate cost, and they are not strongly redundant with
another retained indicator \parencite{andrews2004,bunemann2018}. Six indicators
met the criteria: soil organic carbon (SOC), soil pH, extractable phosphorus
(P), bulk density (BD), water-stable aggregates (WSA) and microbial biomass
carbon (MBC). Together they cover chemical, physical and biological properties
\parencite{doran2000}. Their laboratory methods are listed in
Table~\ref{tab:scoring}; SOC was determined by wet oxidation
\parencite{walkley1934} and P with the Mehlich~3 extractant
\parencite{mehlich1984}.

\input{\subfix{../tables/scoring}}

\subsection{Scoring}
Each measured value $x$ was converted to a score $s$ between 0 and 1. For an
indicator where more is better,
\begin{equation}
  s = \min\!\left(1,\ \max\!\left(0,\ \frac{x - L}{U - L}\right)\right),
  \label{eq:more}
\end{equation}
and for an indicator where less is better the score is $1 - s$ with the same
limits. For soil pH, which has an optimum range from $O_1$ to $O_2$, the score
is 1 inside the range and falls linearly to 0 at $L$ below and $U$ above:
\begin{equation}
  s =
  \begin{cases}
    0 & x \le L \text{ or } x \ge U,\\[2pt]
    \dfrac{x - L}{O_1 - L} & L < x < O_1,\\[6pt]
    1 & O_1 \le x \le O_2,\\[2pt]
    \dfrac{U - x}{U - O_2} & O_2 < x < U.
  \end{cases}
  \label{eq:optimum}
\end{equation}
The limits in Table~\ref{tab:scoring} follow the shape of the scoring curves of
the Soil Management Assessment Framework \parencite{andrews2004}; their values
were set for this volume and are not calibrated for any real soil.

\subsection{Soil health index}
The soil health index (SHI) of a plot is the mean of its six indicator scores,
\begin{equation}
  \mathrm{SHI} = \frac{1}{n}\sum_{i=1}^{n} s_i, \qquad n = 6,
  \label{eq:shi}
\end{equation}
so that every indicator carries equal weight and the index lies between 0 and
1. The index of each system was compared between systems within each site.
